La gesti\'on eficiente del capital constituye uno de los problemas
fundamentales de las finanzas modernas. El problema de optimizaci\'on de
portafolios, formalizado por Markowitz (1952), establece los fundamentos
matem\'aticos de la diversificaci\'on y define la frontera eficiente como el
conjunto de portafolios \'optimos en el espacio riesgo-retorno. En la
pr\'actica, el universo de activos puede comprender cientos de instrumentos,
los clientes presentan perfiles de riesgo heterog\'eneos y las condiciones del
mercado evolucionan continuamente. En este contexto, la consolidaci\'on de
plataformas de asesor\'ia financiera automatizada ha generado una demanda
creciente de sistemas capaces de producir recomendaciones personalizadas,
fundamentadas y actualizadas de forma peri\'odica \cite{dacunto2019}.

El presente informe documenta la segunda fase del proyecto FinPUC, una
plataforma ficticia de recomendaci\'on de inversiones desarrollada en el marco
del curso Taller de Investigaci\'on Operativa (ICS2122) de la Pontificia
Universidad Cat\'olica de Chile. Se presenta la implementaci\'on computacional completa de dicho modelo,
incluyendo tres capas metodol\'ogicas: Markowitz, Black-Litterman con cuatro variantes de \textit{views} y simulaci\'on Monte Carlo; resultados
comparativos contra un caso base equiponderado, y validaci\'on mediante
m\'ultiples horizontes de entrenamiento y escenarios con y sin pandemia.

El sistema opera en dos escalas temporales complementarias. A nivel de
optimizaci\'on, los pesos del portafolio se recalculan cada semestre (126
d\'ias h\'abiles bursatil) mediante rebalanceo con estimaci\'on actualizada de
par\'ametros. A nivel de simulaci\'on del cliente, el modelo Monte Carlo
eval\'ua semanalmente durante 260 semanas si el cliente acepta la
recomendaci\'on vigente o retira sus fondos, generando m\'etricas de riqueza
terminal, probabilidad de ganancia y utilidad de la empresa. La decisi\'on
central consiste en determinar la fracci\'on del capital a asignar a cada
acci\'on de un universo de 601 activos, maximizando simult\'aneamente el
retorno esperado del inversionista y la utilidad de FinPUC.